This paper introduced KCH-MedRank, a knowledge-constrained learning-to-rank framework for biomedical evidence reranking. The central finding is that flat biomedical knowledge features explain most aggregate ranking gains. Hypergraph structure provides targeted but modest early-rank improvements under indirect-evidence conditions---specifically, questions lacking direct MeSH overlap and multi-evidence scenarios---but does not significantly exceed the flat knowledge baseline at the standard top-10 evaluation depth.

On BioASQ, KCH-MedRank significantly outperforms a true MedCPT Cross-Encoder baseline in Recall@10 ($+0.0157$, $p=0.0076$) and runs $18.6\times$ faster at reranking time (reranking-stage latency under cached features, excluding candidate generation). Flat biomedical knowledge features carry most of the aggregate gains. The hypergraph contributes targeted early-rank benefit at $k{=}5$, concentrated in questions that lack direct MeSH overlap between query and evidence. Controlled ablation suggests that removing domain-grounded concept annotations degrades hypergraph performance; however, a full $2 \times 2$ factorial design is needed to cleanly separate the effects of graph topology from the effects of medical knowledge features. BEIR NFCorpus confirms that supervised retrieval-feature reranking generalizes, and the PubMedQA diagnostic suggests that evidence ordering affects downstream answer selection under controlled conditions.

The current evidence supports the conclusion that medical knowledge constraints are necessary for productive graph-based biomedical reranking, and that flat biomedical knowledge features are the primary driver of improvement. The ``raw topology $\to$ noise, medical constraints $\to$ signal'' transformation narrative requires the $2 \times 2$ factorial ablation noted above. An interpretability analysis categorizes 648 of 661 rescued gold passages by dominant activated knowledge mechanism; MeSH hierarchy paths account for 75.9\% of rescues. A lightweight synonym-aware normalization module, integrated into the feature pipeline, bridges 18.6\% of the zero-overlap failure cases identified in error analysis, yielding additional gains of Recall@10 $+0.0056$ and nDCG@10 $+0.0090$ on the validation split. The remaining synonym gap---cases where MeSH entry terms do not cover ad-hoc biochemical nomenclature---points to UMLS concept normalization as a natural extension.

KCH-MedRank is an efficient, interpretable biomedical evidence reranking module. The code and processed feature files will be released upon publication.
